\documentclass[11pt]{article}
\usepackage{tcee}
\addbibresource{ejemplo.bib}
\tcetema{3.A.0}{Tema de ejemplo. Equilibrio parcial y estática comparativa}
\tcefecha{10/10/2026}
\begin{document}
\tceetitulo
\section*{Introducción}
\begin{itemize}
\item \textbf{\emph{Enganche}}: \autor{Alfred Marshall}, en sus \emph{Principios de Economía} (1890), define la economía como \emph{la ciencia de la vida diaria}.\footcite[cap.~3]{MasColell1995}
  \begin{itemize}
  \item El \concepto{individualismo metodológico} \vertema{3.A.16}.
    \begin{itemize}\item Tercer nivel.\end{itemize}
  \end{itemize}
\end{itemize}
\begin{notaopositor}
En el cante basta con enunciar los axiomas y dibujar el gráfico 1.
\end{notaopositor}
\section{Equilibrio}
\subsection{Definición}
El equilibrio cumple $D(P^*)=S(P^*)$:
\begin{equation}
  Q^* = D(P^*) = S(P^*) \label{eq:equilibrio}
\end{equation}
La tasa de paro fue del 10,3\,\%.\footcite{INE2026EPA}
\begin{figure}[H]
\centering
\caption{Equilibrio y desplazamiento de la demanda}\label{fig:equilibrio}
\grafico[0.55\textwidth]{oferta-demanda}
\fuente{elaboración propia a partir de \textcite{MasColell1995}.}
\end{figure}
\begin{ideaclave}
Un aumento de la demanda eleva el precio y la cantidad de equilibrio.
\end{ideaclave}
\subsubsection{Un nivel más}
Texto.
\begin{anotaciones}
Lectura recomendada: \url{https://www.revistasice.com/}
\end{anotaciones}
\section*{Conclusión}
Cierre.
\bibliografia
\end{document}
